% ============================================================================
% Chapitre 6 — Réalisation et implémentation
% ============================================================================
\chapter{Réalisation et implémentation}

\section*{Introduction}

Ce chapitre présente les détails d'implémentation des modules clés de \caimp{},
illustrés par des extraits de code commentés. Nous décrivons les difficultés
rencontrées, les solutions adoptées, et la répartition du travail au sein
de l'équipe.

\section{Organisation du dépôt}

Le dépôt \caimp{} v2 est organisé selon une structure monorepo, chaque service
occupant son propre répertoire sous \texttt{services/} :

\begin{lstlisting}[language=YAML, caption={Structure du dépôt CAIMP}]
caimp/
├── services/
│   ├── agent/           # Python : collecte + enrôlement mTLS
│   ├── telemetry-writer/# Go : ingestion OTLP + détection
│   ├── ai-worker/       # Python : NATS → RAG → Ollama
│   ├── alert-engine/    # Python : règles + déduplication
│   ├── forecast-engine/ # Python : Holt-Winters + LLM
│   ├── correlation-engine/ # Python : corrélation changements
│   ├── change-tracker/  # Python : webhooks + ingest
│   ├── api-admin/       # Python FastAPI : auth + PKI
│   ├── api-query/       # Python FastAPI : lecture + PDF
│   ├── ws-gateway/      # Python : NATS → WebSocket
│   └── ai-orchestrator/ # Python FastAPI : Splunk + Chat
├── infra/
│   ├── postgres/migrations/ # 14 migrations SQL
│   ├── nats/            # config + setup_streams.py
│   └── nginx/           # nginx.conf + certs
├── frontend/            # React 18 + Vite + TypeScript
├── libs/nats_client/    # Client NATS partagé (Python)
├── scripts/             # seed_demo.sql, smoke_test
├── docker-compose.yml
└── Makefile
\end{lstlisting}

\section{Répartition des tâches}

\begin{table}[H]
\centering
\caption{Répartition des responsabilités entre les membres de l'équipe}
\label{tab:repartition}
\begin{tabularx}{\textwidth}{lXl}
\toprule
\textbf{Membre} & \textbf{Modules principaux} & \textbf{Semaines} \\
\midrule
CHAHIDY Rim  & Telemetry Writer (Go), Alert Engine, frontend React
               (Dashboard, Forecasts, Incidents), CSS/design system & S1--S4 \\
SAMDI Wiam   & AI Worker, AI Orchestrator, Correlation Engine,
               RAG pipeline, prompts LLM, PDF reports & S1--S4 \\
TAÏK Youssef & Admin API, Query API, base de données (migrations RLS),
               Change Tracker, agent CAIMP, DevOps (Docker Compose,
               Makefile, scripts) & S1--S4 \\
\midrule
\multicolumn{2}{l}{Travaux communs : architecture, revues de code, tests d'intégration,
  documentation} & S1--S4 \\
\bottomrule
\end{tabularx}
\end{table}

\section{Module 1 : Agent de supervision (Python)}

\subsection{Collecte de métriques}

L'agent utilise la bibliothèque \texttt{psutil} pour lire les métriques système
de manière portable, puis les exporte via l'OpenTelemetry SDK :

\begin{lstlisting}[language=Python, caption={Extrait de \texttt{caimp\_agent/collectors.py}}]
from opentelemetry import metrics

def register_all(meter: metrics.Meter) -> None:
    """Enregistre les collecteurs CPU, mémoire et disque."""

    # CPU : utilisation en pourcentage (0.0-1.0)
    cpu = meter.create_observable_gauge(
        "cpu_percent",
        callbacks=[lambda: [(psutil.cpu_percent(interval=1) / 100.0, {})]],
        description="CPU utilization",
        unit="1",
    )

    # Mémoire virtuelle
    mem = meter.create_observable_gauge(
        "memory_percent",
        callbacks=[lambda: [(psutil.virtual_memory().percent / 100.0, {})]],
        description="Memory utilization",
        unit="1",
    )

    # Disque (partition principale)
    def _disk():
        try:
            usage = psutil.disk_usage("/")
            return [(usage.percent / 100.0, {})]
        except PermissionError:
            return [(0.0, {})]

    disk = meter.create_observable_gauge(
        "disk_percent",
        callbacks=[_disk],
        description="Disk utilization",
        unit="1",
    )
\end{lstlisting}

\subsection{Enrôlement mTLS}

L'enrôlement est la séquence critique de sécurité. L'agent génère une paire de
clés ECDSA P-256 localement, transmet la clé publique à l'Admin API, qui l'emballe
dans un certificat signé par la CA interne.

\section{Module 2 : Telemetry Writer (Go)}

\subsection{Détecteur d'anomalies}

Le détecteur Go implémente les trois algorithmes de manière concurrente. Voici
l'implémentation du détecteur Z-score :

\begin{lstlisting}[language=C, caption={Extrait du détecteur Z-score (\texttt{detector.go})}]
// zScore calcule le score Z sur l'historique des 100 dernières valeurs.
func (d *Detector) zScore(
    ctx context.Context,
    points []receiver.MetricPoint,
) []AnomalyEvent {
    var events []AnomalyEvent
    for _, p := range points {
        // Récupère l'historique depuis PostgreSQL
        rows, err := d.pool.Query(ctx, `
            SELECT value FROM metrics
            WHERE server_id=$1 AND metric_name=$2
            ORDER BY time DESC LIMIT 100`, p.ServerID, p.MetricName)
        if err != nil {
            continue
        }
        var vals []float64
        for rows.Next() {
            var v float64
            if err := rows.Scan(&v); err == nil {
                vals = append(vals, v)
            }
        }
        rows.Close()
        if len(vals) < 30 { // pas assez d'historique
            continue
        }
        // Calcul μ et σ
        mean, std := meanStd(vals)
        if std < 1e-9 {
            continue
        }
        z := math.Abs((p.Value - mean) / std)
        if z >= d.cfg.ZScoreThreshold {
            sev := "warning"
            if z >= d.cfg.ZScoreThreshold*1.5 {
                sev = "critical"
            }
            events = append(events, AnomalyEvent{
                IncidentID: uuid.New().String(),
                OrgID:      p.OrgID,
                ServerID:   p.ServerID,
                MetricName: p.MetricName,
                Value:      p.Value,
                Detector:   "zscore",
                Severity:   sev,
                Timestamp:  p.Time,
            })
        }
    }
    return events
}
\end{lstlisting}

\section{Module 3 : AI Worker et pipeline RAG}

\subsection{Construction du prompt}

Le prompt est construit en plusieurs étapes : récupération du contexte historique,
recherche RAG dans pgvector, et assemblage.

\begin{lstlisting}[language=Python, caption={Extrait de \texttt{prompt\_builder.py}}]
def build_prompt(event: dict, metric_history: list[dict],
                 rag_docs: list[dict]) -> str:
    """
    Construit un prompt ciblant un développeur sans formation DevOps.
    Les métriques sont affichées en pourcentages lisibles.
    """
    history_lines = "\n".join(
        f"  {r['time']}: {r['value']:.1%}"
        if r['value'] <= 1.0 else f"  {r['time']}: {r['value']:.1f}%"
        for r in metric_history[:20]
    ) or "  Pas d'historique disponible"

    rag_lines = "\n\n".join(
        f"[{d['type'].upper()}] {d['title'] or 'Incident passé'}:\n"
        f"{d['content'][:400]}"
        for d in rag_docs
    ) or "  Aucun incident similaire en base"

    pct = f"{event.get('value', 0) * 100:.1f}%"
    return f"""Tu aides un développeur sans formation DevOps...
ALERTE : {event.get('metric_name')} a atteint {pct}
HISTORIQUE : {history_lines}
INCIDENTS SIMILAIRES : {rag_lines}
{_OUTPUT_SCHEMA}"""
\end{lstlisting}

\section{Module 4 : Moteur de prévision (Holt-Winters)}

\subsection{Implémentation NumPy}

L'implémentation du lissage exponentiel double est entièrement vectorisée
pour minimiser la latence :

\begin{lstlisting}[language=Python, caption={Extrait de \texttt{forecaster.py} --- lissage de Holt-Winters}]
def _holt_linear(
    values: np.ndarray, horizon: int,
    alpha: float = 0.3, beta: float = 0.1
) -> tuple[np.ndarray, float, np.ndarray]:
    """Lissage exponentiel double. Retourne (prévisions, tendance, ajusté)."""
    n = len(values)
    # Initialisation
    level = float(values[0])
    trend = float((values[-1] - values[0]) / max(n - 1, 1))
    fitted = np.empty(n)

    for t in range(n):
        fitted[t] = level + trend
        if t < n - 1:
            prev_level = level
            level = alpha * float(values[t]) + (1.0 - alpha) * (level + trend)
            trend = beta * (level - prev_level) + (1.0 - beta) * trend

    # Prévisions : L_t + h * T_t
    forecasts = np.array([level + (h + 1) * trend for h in range(horizon)])
    return forecasts, trend, fitted
\end{lstlisting}

\section{Module 5 : Moteur de corrélation de changements}

La fonction de scoring implémente la formule $s = e^{-0.1 \Delta t} \times w_{\text{type}}$ :

\begin{lstlisting}[language=Python, caption={Fonction de scoring dans \texttt{correlation-engine/main.py}}]
# Poids par type de changement (plus = plus susceptible de causer une anomalie)
_TYPE_WEIGHT = {
    "deployment":  1.0,
    "config":      0.9,
    "package":     0.8,
    "docker_pull": 0.75,
    "restart":     0.7,
    "cron":        0.5,
    "ssh_login":   0.3,
}

def _score_change(change: dict, anomaly_time: datetime) -> float:
    """Calcule le score de corrélation d'un changement avec une anomalie."""
    occurred_at = change["occurred_at"]
    if occurred_at.tzinfo is None:
        occurred_at = occurred_at.replace(tzinfo=timezone.utc)

    # Délai en minutes entre le changement et l'anomalie
    minutes_before = (anomaly_time - occurred_at).total_seconds() / 60
    if minutes_before < 0:
        return 0.0  # Changement postérieur à l'anomalie

    # Décroissance exponentielle : f(dt) = exp(-lambda * dt)
    temporal = math.exp(-0.1 * max(0, minutes_before))
    type_weight = _TYPE_WEIGHT.get(change["change_type"], 0.5)
    return temporal * type_weight
\end{lstlisting}

\section{Module 6 : Frontend React}

\subsection{Tableau de bord orienté réponses}

Le composant central du nouveau tableau de bord est \texttt{StatusBanner},
qui affiche immédiatement le statut global sans nécessiter de lecture de graphes :

\begin{lstlisting}[language=C, caption={Extrait de \texttt{Dashboard.tsx} --- composant StatusBanner}]
function StatusBanner({ summary }: { summary: DashboardSummary }) {
  const isCritical = summary.overall_status === 'critical'
  const isHealthy  = summary.overall_status === 'healthy'
  const bg = isCritical ? 'var(--danger)' :
             isHealthy  ? 'var(--success)' : 'var(--warning)'

  return (
    <div style={{ background: faint, border: `1.5px solid ${bg}`,
                  borderRadius: 14, padding: '20px 28px',
                  display: 'flex', alignItems: 'center', gap: 16 }}>
      {/* Icône de statut */}
      <div style={{ width: 52, height: 52, borderRadius: '50%',
                    background: bg, display: 'flex',
                    alignItems: 'center', justifyContent: 'center' }}>
        {isCritical ? <AlertTriangle size={28} color="#fff" />
                    : <Check size={28} color="#fff" />}
      </div>
      {/* Message principal */}
      <div>
        <div style={{ fontSize: 18, fontWeight: 800, color: bg }}>
          {isHealthy ? 'Tous les systèmes fonctionnent' :
           isCritical ? 'Action requise' :
           'Problèmes potentiels détectés'}
        </div>
        <div style={{ fontSize: 13.5, color: 'var(--text-2)' }}>
          {summary.overall_message}
        </div>
      </div>
    </div>
  )
}
\end{lstlisting}

\section{Module 7 : Génération de rapports PDF}

Le rapport PDF est généré avec ReportLab Platypus. La page de couverture
est dessinée directement sur le canvas de la première page :

\begin{lstlisting}[language=Python, caption={Extrait de \texttt{reports.py} --- page de couverture}]
def _cover_chrome(canvas, doc):
    canvas.saveState()
    # Fond brun foncé
    canvas.setFillColor(colors.HexColor("#3d2b1f"))
    canvas.rect(0, 0, PW, PH, fill=1, stroke=0)

    # Cercles décoratifs
    canvas.setFillColor(colors.HexColor("#5c4033"))
    canvas.circle(PW + 8*mm, PH - 18*mm, 88*mm, fill=1, stroke=0)

    # Logo "C" en badge arrondi
    canvas.setFillColor(WHITE)
    canvas.roundRect(MARGIN, PH*0.60, 22*mm, 22*mm, 3.5*mm, fill=1, stroke=0)
    canvas.setFillColor(colors.HexColor("#3d2b1f"))
    canvas.setFont("Times-Bold", 17)
    canvas.drawCentredString(MARGIN + 11*mm, PH*0.60 + 6.5*mm, "C")

    # Titre du rapport
    canvas.setFont("Helvetica-Bold", 34)
    canvas.setFillColor(WHITE)
    canvas.drawString(MARGIN, PH*0.475, "Infrastructure")
    canvas.drawString(MARGIN, PH*0.475 - 14*mm, "Health Report")
    canvas.restoreState()
\end{lstlisting}

\section{Difficultés rencontrées et solutions}

\begin{table}[H]
\centering
\caption{Principales difficultés rencontrées et solutions apportées}
\label{tab:difficultes}
\small
\begin{tabularx}{\textwidth}{lXX}
\toprule
\textbf{Difficulté} & \textbf{Impact} & \textbf{Solution} \\
\midrule
Désynchro des noms de métriques OTLP
  & L'agent utilisait \texttt{cpu\_percent}, le dashboard cherchait
    \texttt{system.cpu.utilization}
  & Table de correspondance dans le Query API \\
Corruption de contexte dans le chat SSE
  & Les tokens SSE arrivaient dans le mauvais ordre
  & Numérotation des frames, buffer ordonné côté client \\
Lenteur des requêtes pgvector sur INDEX IVFFlat
  & Recherche RAG > 500ms avec 10k vecteurs
  & Migration vers index HNSW, 8x plus rapide \\
Mémoire insuffisante pour Llama 3.1 8B sur 4 Go RAM
  & OOM killer termine Ollama
  & Mode test : switch vers llama3.2:1b, 4x plus petit \\
Déduplication insuffisante des alertes
  & 50+ alertes identiques en 15 min après un pic CPU
  & Clés NX Redis avec TTL 15 min + détection de burst \\
\bottomrule
\end{tabularx}
\end{table}

\section*{Conclusion}

Ce chapitre a exposé les détails d'implémentation des sept modules principaux de
\caimp{}, illustrant la diversité des langages et des paradigmes mobilisés.
Les extraits de code montrent comment les principes architecturaux du chapitre 4
se traduisent en lignes de code concrètes. Le chapitre suivant évalue la
correctness de ces implémentations par une batterie de tests systématiques.
